\paragraph{Circumferential fibers.} The coupled 3D cylinder reproduces the PV loop of the
0D block in the same circulation. Volume, stroke volume, ejection fraction, wall
displacement and strains converge with the mesh (volume and stroke volume at about third
order; the volume error falls to $4\cdot10^{-4}$ at nr = 4). The cavity pressure, the
valve flows and $\dot P$ stop improving from nr = 2 on. Halving the time step (800 instead
of 400 steps per cycle, in both models) halves this level: the pressure error at nr = 4
goes from $2.3\cdot10^{-3}$ to $8.8\cdot10^{-4}$, the inflow error from $1.4\cdot10^{-2}$
to $3.7\cdot10^{-3}$, and the valve times differ by less than one time step. The
remaining pressure error is therefore the $O(\Delta t)$ difference between the coupled 3D
(generalized-$\alpha$) and the stand-alone 0D time integration, not a model difference.
The second cycle, in which the 0D model is periodic, gives the same errors as the first.
At nr = 8 (19584 elements, on the cluster) the stroke volume and the ejection fraction
converge further at order 3.3 to $4\cdot10^{-5}$, and the strains at 1.2--2.1, while the
maxima of the volume, pressure and $\rho$ differences stay at their nr = 4 values
($3.9\cdot10^{-4}$, $2.3\cdot10^{-3}$, $1.5\cdot10^{-3}$). All three now peak at the start of
the compared cycle ($t = 0.02$~s). Split into windows of the cycle, the volume and $\rho$
differences in the ejection (0.2--0.35~s) fall with the mesh (volume $1.6\cdot10^{-4}$ to
$6.3\cdot10^{-5}$, $\rho$ $1.2\cdot10^{-3}$ to $2.5\cdot10^{-4}$ from nr = 4 to 8) and not with
the time step; early in the cycle (0.02--0.2~s) they fall with the time step (volume
$3.9\cdot10^{-4}$ to $1.9\cdot10^{-4}$) and not with the mesh. The 3D model thus converges in
space to the 0D block down to the time-integration difference.

\paragraph{Thick wall.} With $R_e/R_i = 2.5$ instead of 1.74 (the same inner radius,
material, tension and circulation) the 0D tube ejects 37\% instead of 27\% of its volume
at a peak pressure of 155~mmHg. The coupled 3D model converges in the same way: volume,
stroke volume, ejection fraction and stroke work at orders 2.6--3.3 down to
$7\cdot10^{-4}$--$1.2\cdot10^{-3}$ at nr = 8, the wall displacement at 2.4 to
$2.4\cdot10^{-3}$. The pressure difference is $2.4\cdot10^{-3}$ at nr = 4 and
$2.3\cdot10^{-3}$ at nr = 8, but for different reasons. At nr = 4 its maximum lies in the
ejection ($t = 0.34$~s) and stays at $2.3\cdot10^{-3}$ with half the time step
(\texttt{active\_thick\_dt2}): there it is the mesh. At nr = 8 the ejection part has
fallen to $5\cdot10^{-4}$--$1.7\cdot10^{-3}$, and the maximum sits at the start of the
compared cycle ($t = 0.02$~s), where half the time step reduces it from $2.3\cdot10^{-3}$ to
$8.6\cdot10^{-4}$ (at nr = 4), as in the base geometry ($2.3\cdot10^{-3}$ to
$8.8\cdot10^{-4}$). Half the time step also reduces the inflow difference
($1.8\cdot10^{-2}$ to $1.1\cdot10^{-2}$), the peak-$\dot P$ difference ($2.4\cdot10^{-3}$ to
$1.0\cdot10^{-3}$) and the valve-time offset (1.5 to 0.7~ms) at nr = 4.

\paragraph{Helical fibers.} With the $\pm60^\circ$ helix of the test case, in the $2L_0$ tube
(the same circulation with twice the initial volume; ejection fraction 0.18, peak
157~mmHg), the contraction twists the 0D tube by up to $15.6^\circ$ over its length, about
90 times the passive twist. Volume, pressure and flows are those of the whole 3D tube,
which drives the circulation; the wall kinematics, strains and the twist rate are taken
at mid-length. From nr = 1 to 2 every difference falls: the volume from $1.7\cdot10^{-2}$ to
$5.5\cdot10^{-3}$, the pressure from $5.7\cdot10^{-2}$ to $1.2\cdot10^{-2}$, the mid-length
$\rho$ from $1.1\cdot10^{-1}$ to $2.5\cdot10^{-2}$, the twist rate from $4.6\cdot10^{-2}$ to
$6.9\cdot10^{-3}$ of its peak, the valve times from 10.6 to 0.9~ms. The stroke volume falls
less ($4.5\cdot10^{-2}$ to $2.6\cdot10^{-2}$): the circulation sees the volume of the whole
tube, end layers included. The 0D replay at the 3D pressure (Section~\ref{sec:replay})
separates the two: at nr = 2 the stroke volume of the whole tube differs from the 0D by
$1.5\cdot10^{-2}$, that of the mid-length section by $9\cdot10^{-4}$. In the original tube
the end layers are twice the share of the volume, and the pressure and stroke-volume
differences stay at 3--6\% from nr = 2 to 4; this is why the coupled case, like the passive
one, uses the $2L_0$ tube.
